\chapter{Testing}\label{app:testing}
\sloppy

Both implementations come with tests. This appendix describes each
suite: what it tests, how to run it, how long it takes, how to run a
single case, and how to add one. All commands are run from the root of
the source tree.

\begin{notebox}
The expected results of the tests come from the 1992 manual, not from
either implementation. Where the manual gives an example with its
result, the test expects that result; where it specifies behavior
without an example, the test is written from the specification and cites
the chapter and section. Neither implementation is taken as the
reference for the other.
\end{notebox}

\section{The suites at a glance}\label{sec:test-overview}

\begin{center}
\small
\begin{tabular}{>{\raggedright\arraybackslash}p{0.33\textwidth}%
                >{\raggedright\arraybackslash}p{0.58\textwidth}}
\toprule
Command & What it tests \\
\midrule
\cmd{sbcl --script tests/front-end/run.lisp} & the compiler's front end:
  reading, modules, conversion to core forms, diagnostics \\
\cmd{sbcl --script tests/ir/run.lisp} & the IR and its analysis \\
\cmd{sbcl --script tests/codegen/run.lisp} & the compiler's features,
  phase by phase: programs compiled, run, and their output compared \\
\cmd{sbcl --script tests/conformance/run.lisp} & the manual's examples
  and specifications, chapter by chapter \\
\cmd{sbcl --script tests/stress/run.lisp} & 10\textsuperscript{8} tail
  calls; many threads at once \\
\cmd{sh tests/driver/run.sh} & the command line, \cmd{bin/dylanc} \\
\cmd{make -C runtime test} & the C runtime's data representation \\
\code{(dylan:test-examples)} & the interactive implementation, on the
  book's examples \\
\cmd{sbcl --script tests/interpreter/run.lisp} & the compiler's test
  programs, run in the interactive implementation \\
\bottomrule
\end{tabular}
\end{center}

Each command prints one line per failure, beginning \code{FAIL}, and a
summary; the suites of the compiler exit with status 0 when everything
passed and 1 otherwise.

\section{The front end and the IR}\label{sec:test-front-ir}

\file{tests/front-end/cases.lisp} holds 55 cases. Each is a list

\begin{lisp}
(name "dylan source" expected)
\end{lisp}

\noindent where \meta{expected} is either the core forms the source
converts to, as the compiler's \code{unparse} writes them, one per line,
or \code{(:diagnostics "text" ...)}: the diagnostics the source must
produce, in order, each given as a string it must contain. For example:

\begin{lisp}
 (no-complex
  "#c(1 2)"
  (:diagnostics "error: complex numbers are not supported"))
\end{lisp}

After the cases, the runner converts \file{dylan-examples.dyl}, the
book's examples, which must produce no diagnostic but the one error that
its page~68 example is meant to cause.

\file{tests/ir/cases.lisp} holds 18 cases of the form

\begin{lisp}
(name "dylan source" "expected IR" check ...)
\end{lisp}

\noindent The expected IR is the A-normal form of each top-level form;
each \meta{check} names a variable or function and properties it must
have --- whether a variable is captured or assigned, how many references
it has, whether a function escapes, where it is called from and which of
those calls are tail calls, and so on (the file's header lists them).
The IR of \file{dylan-examples.dyl} must then pass the checker of the
IR's invariants.

Neither suite runs a C compiler. Given the argument
\cmd{print}, each shows what every case converts to, which is how to
write the expected text of a new case:

\begin{shell}
$ sbcl --script tests/front-end/run.lisp print
\end{shell}

\section{Programs compiled and run}\label{sec:test-programs}

Three suites --- \file{tests/codegen}, \file{tests/conformance} and
\file{tests/stress} --- consist of whole Dylan programs that are
compiled, built, run, and whose output is compared with what the test
expects. They share one runner, \file{tests/codegen/runner.lisp}.

\begin{description}
\item[\file{tests/codegen/cases*.lisp}] 108 programs testing the
  compiler's features phase by phase: \file{cases.lisp} (calls, tail
  calls, closures, \code{\#rest} and \code{\#key}, multiple values,
  run-time errors), \file{cases-control.lisp} (exits, cleanups,
  conditions), \file{cases-objects.lisp} (classes, slots, dispatch,
  \code{next-method}, reflection), \file{cases-library.lisp},
  \file{cases-threads.lisp}, \file{cases-macros.lisp} and
  \file{cases-optimize.lisp} (type inference, inlining, unboxing,
  escape analysis).
\item[\file{tests/conformance/ch*.lisp}] 237 programs, one file for each
  chapter of the manual from 2 to 18: the manual's examples with the
  results it gives, and its specifications, each test citing its
  section. What a static program cannot do is not tested here; it is
  listed in Appendix~\ref{app:differences}.
\item[\file{tests/stress/cases-*.lisp}] 11 programs that take minutes:
  $10^8$ tail calls through unknown functions, generic functions,
  closures, \code{next-method} and \code{apply}, each in a 256~KB stack;
  and many threads allocating, contending for a bounded buffer,
  interning symbols, signalling conditions, binding thread variables
  and making threads.
\end{description}

\subsection{The case format}\label{sec:test-case-format}

\index{test cases}
A file of cases holds one list of cases, each

\begin{lisp}
(name "dylan source" "expected standard output"
      [:error "text"] [:compile-error "text"] [:compiler-only t])
\end{lisp}

The program is the source string; the strings follow Lisp's rules, so a
double quote inside is written \code{\textbackslash"} and a backslash
\code{\textbackslash\textbackslash}. The output must be exactly the
expected text, newline for newline. With \code{:error}, the program must
instead stop with status~1 and an error message on the standard error
containing the text --- after writing the expected output. With
\code{:compile-error}, the compiler must report an error containing the
text, and the program is not run. \code{:compiler-only t} marks a case
that tests what the interactive implementation does not do by design ---
63-bit integers, a static program's restrictions, the compiler's own
diagnostics, threads --- which \file{tests/interpreter/run.lisp}
(Section~\ref{sec:test-interactive}) skips. From
\file{cases-threads.lisp}:

\begin{lisp}
 (thread-values
  "(define t1 (make-thread (method () (values 1 2))))
   (bind ((a b (join-thread t1))) (print (list a b)))
   (print (join-thread t1))
   (define t2 (make-thread (method () (current-thread))))
   (print (list (id? (join-thread t2) t2) (id? (current-thread) t2)))
   (print (instance? t2 <thread>))"
  "(1 2)
1
(#t #f)
#t
")
\end{lisp}

\noindent and from \file{cases-macros.lisp}, a case that must fail to
compile:

\begin{lisp}
 (macro-not-top-level
  "(define f (method () (define-macro m () 1)))"
  ""
  :compile-error "define-macro must be at top level")
\end{lisp}

Threads run in any order, so a test of threads must print only what does
not depend on the order. Every test program runs with a 256~KB stack
(\cmd{ulimit -s 256}), so a call in tail position that was not compiled
as a jump overflows it quickly and the test fails.

\subsection{Modes}\label{sec:test-modes}

\index{test modes}
Each program is compiled to C once and then built and run in each
\emph{configuration} of a \emph{mode}, named on the command line:

\begin{description}
\item[\cmd{all}] (the default) GCC and Clang, each with \code{musttail}
  and with trampolines: four builds of each program, the first at
  \cmd{-O2} and the others at \cmd{-O0}.
\item[\cmd{quick}] The first of those only: GCC, \code{musttail},
  \cmd{-O2}.
\item[\cmd{whole}] As \cmd{quick}, but with the library compiled with
  each program and tree shaken (\cmd{-{}-whole}), not compiled apart.
\item[\cmd{gc-stress}] GCC, with \code{musttail} and with trampolines,
  each forcing a collection every 1000 allocations
  (\cmd{-{}-gc-stress} with \code{DYLAN\_GC\_STRESS=1000}).
\item[\cmd{address}] AddressSanitizer and UndefinedBehaviorSanitizer,
  GCC and Clang, in debug builds.
\item[\cmd{thread}] ThreadSanitizer, GCC, in a debug build.
\end{description}

Every mode builds with the default collector, the Memory Pool System.
With the environment variable \code{DYLAN\_GC=boehm} set, it builds with
the Boehm-Demers-Weiser collector instead (Section~\ref{sec:dylanc-gc});
the compiler and runtime should pass every suite with either.

\begin{shell}
$ sbcl --script tests/codegen/run.lisp quick
$ sbcl --script tests/conformance/run.lisp            # all
$ sbcl --script tests/stress/run.lisp thread
\end{shell}

The builds go in \file{runtime/build/tests/} (a subdirectory named for
the mode, except in \cmd{all}). The sanitizer and stress modes build
their own variants of the runtime library the first time they need them.

\subsection{Running some of the cases}\label{sec:test-single}

Further arguments name the case files to run, and the cases: any
argument ending in \file{.lisp} is a file to read instead of the suite's
own, and any other (besides the mode) is the name of a case to run,
whichever file holds it.

\begin{shell}
$ sbcl --script tests/codegen/run.lisp quick thread-values
$ sbcl --script tests/codegen/run.lisp quick tests/codegen/cases-macros.lisp
$ sbcl --script tests/conformance/run.lisp quick ch08-true-false
\end{shell}

The summary line counts the programs, the configurations and the
failures, as in \code{1 program in 1 configuration, 0 failures}.

\subsection{How long they take}\label{sec:test-time}

The programs are built and run one after another, so the time grows
with the number of programs and configurations. A build takes one to two
seconds once the compiler and the library's object are cached
(Section~\ref{sec:dylanc-cache}); a configuration whose C compiler and
options are new must first compile the library, which takes some
seconds more. On the 32-thread machine of Chapter~\ref{ch:performance},
the conformance suite in mode \cmd{all} --- 237 programs in four
configurations --- took about a quarter of an hour; \cmd{quick} needs a
fraction of that. The \cmd{address}, \cmd{thread} and \cmd{gc-stress}
modes are much slower, because their programs run slowly, and the stress
suite takes minutes in any mode. The front-end and IR suites run no C
compiler and are quick. Use \cmd{quick} while working, and \cmd{all} and the other
modes before committing a change to the compiler or runtime.

\section{The command line}\label{sec:test-driver}

\cmd{sh tests/driver/run.sh} tests \cmd{bin/dylanc} itself, in a
temporary directory: the default executable name and \cmd{-o}; that the
C is removed unless \cmd{-k} is given, and what \cmd{-k} keeps; the
\code{\#line} directives and \cmd{-{}-no-lines}; \cmd{-C}, and compiling
its C by hand; each \cmd{-O} level, \cmd{-Os} and \cmd{-g};
\cmd{-{}-trampoline}, \cmd{-{}-cc}, \code{CC}, \cmd{-{}-lto} with GCC and
with Clang, and \cmd{-{}-whole}; the library's object cache, and that two
programs' library C is the same text; a program of several modules,
with a file for each; a macro exported from its module, and \cmd{-E};
run-time errors and their exit status; and the exit status of compile
errors and of bad command lines. It prints \code{FAIL} and a description
for each check that fails, and a count of checks and failures.

\section{The runtime}\label{sec:test-runtime}

\cmd{make -C runtime test} builds and runs
\file{runtime/build/test-runtime}, a C program that tests the runtime's
data representation without any Dylan program: the built-in classes as
the compiler emits them (which \file{tests/runtime/generate.lisp}
writes): immediate values, whether each class is a subclass of each
other, how constants print, integer arithmetic at the edges of the
63-bit range and its errors, other numbers, strings and symbols, the
collector, and threads. It runs in moments once built.

\section{The interactive implementation}\label{sec:test-interactive}

\index{test-examples@\code{test-examples}}
The interactive implementation is tested on the examples from the Dylan
book, \file{dylan-examples.dyl}:

\begin{shell}
$ cd src
$ sbcl --eval '(load "load-dylan.lisp")' \
       --eval '(load "full-test.lisp")' \
       --eval '(dylan:test-examples)'
\end{shell}

\code{test-examples} reads the file form by form, echoes each form,
evaluates it, and prints its values --- or, for a form that signals a
condition, \code{;Error:} and the condition's message, after which it
goes on with the next form. It does not compare the results itself. The
expected results are in \file{src/reference/}, the transcripts Thomas
produced from the same file in 1992 (\file{MIT.output} among them);
compare the output with them. Where they and the manual disagree, the
manual decides; the few remaining differences are listed in
\file{src/README}. Run it from the
\file{src} directory, as shown, since \code{test-examples} finds the
examples at \file{../dylan-examples.dyl}; give it another file name as
its argument to run other examples the same way.

\index{tests/interpreter/run.lisp@\file{tests/interpreter/run.lisp}}
The programs the compiler is tested on run in the interactive
implementation too, by \file{tests/interpreter/run.lisp}, which is how
the two implementations are kept in agreement. It takes files of cases in
the form of Section~\ref{sec:test-case-format}, and test names, as the
compiler's runners do; by default all of \file{tests/codegen/cases*.lisp}.
Each program runs in \code{dylan-user}, as a compiled program does, and
what it defined is forgotten before the next; its output must be the
expected output, and the error it signals, if the case expects one, must
have a message containing the expected text --- whether the compiler
would find the error compiling or running makes no difference here, and
the \code{Dylan error:} of a compiled program's report is not expected.
Cases marked \code{:compiler-only} are skipped and counted:

\begin{shell}
$ sbcl --script tests/interpreter/run.lisp
97 programs, 0 failures, 14 compiler-only skipped
$ sbcl --script tests/interpreter/run.lisp tests/codegen/cases-control.lisp
\end{shell}

Every program not so marked must pass in both implementations.

\section{Adding a test}\label{sec:test-adding}

\begin{enumerate}
\item Choose the suite. A feature of the compiler goes in the
  \file{tests/codegen/cases-*.lisp} file for its area; an example or
  specification from the manual goes in the
  \file{tests/conformance/ch}\meta{nn}\file{.lisp} of its chapter, with
  a comment citing the section and saying whether the expected result
  is the manual's example or follows from its text; a test that takes
  long goes in \file{tests/stress/}.
\item Write the case in the format of Section~\ref{sec:test-case-format}:
  a unique name, the program, and the output the manual or the
  specification requires --- not merely what the implementation prints
  today. Print every result with \code{print} or \code{format}, and,
  for threads, print only what does not depend on their order.
\item Run it alone in \cmd{quick} mode until it passes, then in the
  default mode, and in \cmd{gc-stress}, \cmd{address} or \cmd{thread}
  if it exercises memory or threads:
\begin{shell}
$ sbcl --script tests/codegen/run.lisp quick my-new-case
$ sbcl --script tests/codegen/run.lisp my-new-case
\end{shell}
\item For the front end or the IR, add a case to its \file{cases.lisp},
  using the runner's \cmd{print} argument to see what the source
  converts to, and checking that by hand against the manual before
  making it the expected text.
\end{enumerate}

A test that fails because the implementation is wrong is a bug to fix
in the implementation, not in the test; a test is changed only when it
misread the manual.

\fussy
